\scanpage{130}
\mspage{117}
\uline{Proposition 6.4.} Soit $M$ une catégorie,
et $\Phi\longleftrightarrow\Psi$ un couple $Q$-clos en \deleted{\unsure{la}}
dualité $(\Phi=\Psi^*,\ \Psi=\Phi_*)$. Soit
$u\colon I\longrightarrow I'$ une flèche dans $\mathrm{Cat}$, d'où
suite de foncteurs adjoints

% slanted
\margin{Sous réserve d'existence des $\underrightarrow{\lim}$ et $\underleftarrow{\lim}$ pertinentes dans $M$. \deleted{\unsure{c'est}} pour simplifier $M$ \unsure{stable} pour $\underrightarrow{\lim}$ et $\underleftarrow{\lim}$.}
\diagram{\begin{tikzcd}[ampersand replacement=\&,column sep=huge]
 M(I)\arrow[r,shift left=2.2ex,"{u_!}"]\arrow[r,shift right=2.2ex,"{u_*}"'] \& M(I')\arrow[l,"{u^*}"'].
\end{tikzcd}}

On définit $\Phi(I),\Psi\langle I\rangle,\Phi(I'),\Psi\langle I'\rangle$
comme dans 6.\unsure{2}.5. On a alors :

\textcircled{\scriptsize a}
\[
 \left\{\begin{aligned}
 u_!(\Phi(I))&\subset\Phi(I')\\
 u^*(\unsure{\Psi|\mathscr{T}|T}\langle I'\rangle))&\subset\Psi\langle I'\rangle.
 \end{aligned}\right.\tag*{(6.4.1)}
\]

\textcircled{\scriptsize b} \deleted{D'autre part,} On a l'équivalence
\[
 u^*(\Phi(I'))\subset\Phi(I)\ \Longleftrightarrow\
 u_*(\Psi\langle I\rangle)\subset\Psi\langle I'\rangle.\tag*{(6.4.2)}
\]

\textcircled{\scriptsize c} Les relations équivalentes (6.4.\unsure{2|1}) ont lieu
dans chacun des trois cas suivants :
\deleted{quand $I$ est discrète, ou}

$1^\circ$) $u$ est \unsure{coétale}

$2^\circ$) $I$ est discrète

$3^\circ$) Pour \unsure{tt} $i'\in\operatorname{Ob} I'$, \deleted{\unsure{toute|les}} \unsure{composante|composantes}
\unsure{connexe|connexes} de $I_{/i'}$ admet un objet final.

\uline{Corollaire} 6.5. Soit \deleted{\unsure{le}} $(W,C,F)$ un triple de
Quillen clos \unsure{séparant} dans $M$, et définissons,
pour tt $I$ dans $\mathrm{Cat}$ $C(I)$, $TC(I)$ et $F(I)$, $TF(I)$
comme dans le th. 6.1. a). On a alors
% archive sheet number: 129

